\documentclass[12pt,a4paper]{article}

\usepackage{url}
\sloppy

\title{Book Proposal: \\ 
%\emph{Creating Concurrent Programs}}
Concurrent Programming: \\ Principles, Techniques and Patterns}

\author{Gavin Lowe\footnote{E-mail: \texttt{gavinlowe1@gmail.com}.}}

\begin{document}
\maketitle

\subsection*{Book overview}

This is a book about designing and writing concurrent programs.  My experience
is that other books on the subject present concurrency primitives, but do not
adequately teach techniques whereby these primitives can be used effectively.
The aim here is to present general patterns and techniques that can be reused,
and so help the reader develop a deeper understanding.

The book is aimed at readers who already have experience of writing sequential
programs, but want to learn how to write concurrent programs.  It could be
read either by students or professional programmers.  In Oxford, the course
corresponds to a second-year undergraduate course, although some universities
may prefer to teach the material in a later year.

The book is example-oriented.  Examples that introduce new concurrency
techniques are deliberately simple, to help the reader understand how the
technique works.  But these lead onto larger, more realistic examples, so that
the reader can see how to use the techniques.  The book emphasises various
patterns.  Exercises are used to help reinforce those patterns. 

The book also includes more foundational material, to help the reader
understand underlying principles, what the correctness criteria are, and how
to avoid pitfalls such as deadlocks.

One theme of the book is modularisation, particularly using \emph{concurrent
  objects}.  Operations on such objects should be thread-safe, i.e.,~different
threads do not interfere with one another.  Client code can use that object
without knowing how the concurrency is implemented, and is often very similar
to corresponding code in a sequential program.  This modularisation makes
development easier: it requires local reasoning about a single object at a
time, rather than global reasoning.

Another theme that runs throughout the book is \emph{testing}, a topic that
most other books ignore.  How can we test concurrent programs?  How can we
perform unit testing on a concurrent object?  Often this requires different
techniques from those for sequential programs or objects.

The book makes use of the Scala programming language, together with my own
Scala Concurrency Library (SCL).  These, together, remove much of the
boiler-plate code present in many other languages, so that the programmer can
concentrate on the main ideas.  However, the techniques are readily
translatable to other languages, given suitable concurrency primitives.  An
appendix describes relevant parts of the Scala language; my aim is that
reasonably experienced programmers should need to refer to the appendix only
occasionally, for specific language features.

The book describes message-passing concurrency first, because I think that
this is the most intuitive style of concurrency.  However, this paradigm is
used to introduce other techniques and patterns of concurrent programs.  The
book then progresses onto synchronous parallel programming, and subsequently
describes monitors and semaphores.

In more detail, chapters cover the following topics.
%
\begin{description}
\item[Chapter 1, Introduction] introduces the subject, describing reasons why
  we might want to use concurrency, and gives a brief overview of concurrent
  architectures.  It explains why concurrent programming entails extra
  challenges, by introducing race conditions and problems caused by cache
  inconsistency.

\item[Chapter 2, Message Passing using Channels] introduces the idea of
  message-passing concurrency, and how to define and use channels in SCL{}.  It
  illustrates the technique with a number of small examples, particularly
  examples using fine-grained concurrency: the aim here is to get the reader
  used to thinking concurrently.  It also includes examples of recursive
  parallelism.  It introduces a simple way of testing, by generating random
  input data, and comparing the results of the concurrent program with that
  from a simple sequential program.  It also describes how to reason about the
  order of actions (using Lamport's ``happens-before'' notation).

\item[Chapter 3, Example: Numerical Integration] describes a larger example,
  namely numerical integration based on the trapezium rule.  Different worker
  threads calculate different parts of the integral.  It illustrates the
  technique of a controller distributing tasks to workers, and combining
  together the subresults into the final result (both using message passing).
  It does this in two ways: where each worker is responsible for a fixed part
  of the integral; and the \emph{bag-of-tasks pattern}, where workers perform
  multiple tasks, helping to smooth out variations in speed between workers.
  The chapter talks a bit about tuning, so as to choose values of parameters
  that make the program run quickly; in particular, it demonstrates how a
  badly tuned concurrent program can perform very poorly.  Finally, it talks
  about encapsulation, and the idea of \emph{concurrent objects}.

\item[Chapter 4, Alternation] describes the technique of alternation, where a
  thread can be willing to communicate on one of several channels, whichever
  channel is ready first.  It describes how to achieve this in SCL, and
  illustrates the idea with a number of small examples.  In particular, it
  uses the classic Dining Philosophers example, and illustrates how a deadlock
  can arise.  This leads into a discussion about debugging concurrent
  programs using logging, and also about how to avoid deadlocks in the first
  place. 

\item[Chapter 5, Clients and Servers] introduces the client-server paradigm,
  where clients send requests to a server, which responds.  The paradigm is
  illustrated using a resource server; various techniques are described.  The
  paradigm is also illustrated using a couple of \emph{synchronisation
    objects} (where two or more threads synchronise, possibly exchanging
  data).  Finally, the chapter discusses how to test these synchronisation
  objects, in particular by logging the results of operations, and then
  checking that the log represents a valid execution.

\item[Chapter 6, Interacting Peers] describes various topologies that can be
  used to connect together peers: a centralised topology, a fully connected
  topology, a ring topology, and a heap topology.  These are all illustrated
  on the same simple example, to allow comparison.  Further, a couple of
  distributed algorithms are presented: leadership election in a ring, and
  building a spanning tree in a general graph.  

\item[Chapter 7, Concurrent Datatypes] describes encapsulation inside
  \emph{concurrent datatypes}.  These provide thread-safe implementations of
  standard datatypes such as queues, stacks, sets and mappings.  The
  implementation of the concurrent datatype is responsible for controlling the
  concurrency: but reasoning about correctness is at a local level, so easier
  than global reasoning.  Client code can mostly use the concurrent datatype
  much the same as sequential code would use the corresponding sequential
  datatype.  The chapter presents both total datatypes (where each operation
  can return in every state), and partial datatypes (where operations may
  block until a natural precondition holds).  It presents the standard
  correctness condition for such concurrent datatypes, namely
  \emph{linearization}.  Further, it presents a technique for testing for
  linearization (encapsulated inside SCL, to make the testing
  straightforward).  The paradigm is illustrated by using two concurrent
  datatypes within a concurrent implementation of (near-)breadth-first search.

\item[Chapter 8, Synchronous Parallel Programming] considers programs that
  proceed in rounds, with a synchronisation between all threads---known as a
  \emph{barrier synchronisation}---at the end of each round.  The technique of
  a barrier synchronisation is extended to allow threads to combine values at
  the end of each round, for example to decide whether to terminate.  The
  technique is illustrated with several examples.  Two variants of the pattern
  are considered.  In one variant, each round is split into a subround where
  shared variables are read, and a subround where variables are written, with
  an additional barrier synchronisation between.  The second variant uses two
  copies of the data, one being read and one being written, with their roles
  swapping at the end of each round.

\item[Chapter 9, Using Explicit Locks] considers how to use explicit locks
  (specifically, locks in the SCL library).  Their use is illustrated by the
  implementation of some concurrent objects.  The chapter then considers
  \emph{fine-grained locking} and \emph{sharding}, both techniques that 
  improve scaling as the number of threads grows.

\item[Chapter 10, Conditions and Signals] describes techniques to allow one
  thread to send a signal to another, via a \emph{condition}.  The technique
  is introduced using examples of a partial queue and a bounded partial queue;
  exercises ask students to adapt the ideas to stacks.  The technique is also
  illustrated in implementations of several synchronisation objects: a
  synchronous channel, an exchanger,  a barrier synchronisation, and a readers
  and writers lock.  Finally they are used in a classical synchronisation
  problem, namely the sleeping barber problem. 

\item[Chapter 11, JVM Monitors] describes the concurrency primitives within
  the JVM, providing locking and signalling.  These are used to implement a
  number of concurrent objects (including using fine-grained locking).  They
  are compared with corresponding primitives within SCL.

\item[Chapter 12, Semaphores] describes how to use semaphores in concurrent
  code, both for mutual exclusion and for signalling.  Again, the ideas are
  presented using a number of examples, particularly of concurrent objects.
  The chapter emphasises the technique of \emph{passing the baton}.  It also
  describes counting semaphores.

\item[Chapter 13, Conclusions] gives an overview, and briefly discusses some
  other techniques (atomics, and the related property of lock freedom;
  futures; and competition parallel).  (This chapter is less developed than
  others;  I might yet decide to remove it.) 

\item[Appendix A, A Brief Introduction to Scala] describes enough Scala for a
  reader unfamiliar with the language to understand the rest of the book.  The
  body of the book includes references to this appendix where appropriate.  My
  aim is that most readers will only have to refer occasionally to the
  appendix. 

\item[Appendix B, Experimental Design] gives a brief overview about designing
  tuning experiments.  
\end{description}


%% I might add an additional final chapter, giving an overview of more advanced
%% concurrency techniques, particularly compare-and-set operations, and the
%% properties of lock freedom and wait freedom.  I don't want to treat these
%% topics in full detail, because they are already very well covered by Herlihy
%% and Shavit (discussed below), and because I don't want the book to get much
%% longer.  ...........

The draft contains two additional appendices, containing answers to exercises.
I propose that Appendix~C will be provided online (covering exercises labelled
with ``$\dag$''); and Appendix~D will be available from me, for bona fide
instructors only.


%%%%%%%%%%%%%%%%%%%%%%%%%%%%%%%%%%%%%%%%%%%%%%%%%%%%%%%

\subsection*{Competition}

I discuss here some other books on concurrent programming.
%
Some books concentrate on a particular programming language, typically named
in the title.  My view is that a textbook should teach general techniques,
rather than concentrate on a particular language.  Certainly, a programming
language is very useful in order to make examples concrete; but it should be
straightforward for the reader to translate those techniques to other
languages (given suitable concurrency primitives).

\emph{Concurrent Programming in Java: Design Principles and Patterns}, Doug
Lee, Addison Wesley.
%\url{https://www.anna.ase.ro/Java/resurse/Concurrent_Programming_in_Java.pdf}
%
%I  think the book is badly written.  
I find that this book does not explain the ideas well, and is rather
confusing.  It is over-dense, and feels like a ``brain-dump''.
%% There are a number of technical issues that I disagree with (or, at least,
%% are likely to be misunderstood).
%The main ideas do not come across well: 
In a number of places, I struggled to
understand what the author was trying to say; somebody who is not already
familiar with the material would surely struggle more.
%
Monitors and semaphores are described; but I don't think there are enough
examples or enough explanation for readers to understand the patterns, and so
be able to write similar programs. 
%
Channels are mentioned, but not used in meaningful examples (the one
half-meaningful example is the dining philosophers).  
% Sec 4.5.1
Synchronous data parallel programming is described briefly, with a single
example. 
% Sec 4.4.3
Concurrent objects are used; but there doesn't seem to be sufficient
discussion of what correctness means; there is no mention of linearization.
% 
It also does not seem to cover the following topics from the proposed book:
systematic testing; tuning; logging for debugging and testing; synchronisation
objects (beyond channels and barriers); sharding.

%%%%%

\emph{Learn Concurrent Programming with Go},  James Cutajar, Manning.
% https://github.com/ytx-readings/go-books/blob/main/books/Learn%20Concurrent%20Programming%20with%20Go%20(James%20Cutajar)%20(Z-Library).pdf
In my opinion, this book covers the material rather slowly (although is mostly
clear); as a result, it covers much less ground than the proposed book.  It
describes various concurrency primitives, and illustrates them with toy
examples; but it's lacking the bigger picture.  I don't think there are enough
examples for the reader to spot the patterns.
%
In particular, it seems not to cover the
following topics from the proposed book: systematic testing; 
%% workers and controller pattern, including the bag-of-tasks pattern;
tuning; encapsulation; logging for debugging and testing; synchronisation
objects (beyond channels and barriers); different topologies for
communications; concurrent datatypes; linearization; synchronous parallel
programming; fine-grained locking; sharding.

%%%%%

\emph{The Art of Multiprocessor Programming}, Maurice Herlihy and Nir Shavit,
Morgan Kaufmann.  
%
This is a more advanced book than the proposed book.  As one point of
comparison, in Oxford, the proposed book corresponds to a second-year
undergraduate course, whereas Herlihy and Shavit's book corresponds to a
fourth-year/master's level course.  In common with the proposed book, it
places an emphasis on concurrent objects, with linearizability as the
correctness condition.  It has less material on how to use those concurrent
objects in complete programs.  It does not cover the following topics from the
proposed book: message passing (Chapters 2--6 of the proposed book);
bag-of-tasks pattern; systematic testing; tuning; logging; synchronous
parallel programming; sharding; semaphores.  Instead, the emphasis is on
atomic variables, and lock-free or wait-free operations.  I consider this to
be a very good book: one of the best in Computer Science, and a suitable
follow-on to the proposed book.

%%%%%

\emph{Concurrent Programming: Principles and Practice}, Gregory R. Andrews,
Pearson.  This is a classic book; but it is now rather dated (1991), and also
out of print.  The book uses pseudocode, which makes the material less
concrete, and creates a disconnect from real programs; in particular, it makes
it harder for an instructor to set practical programming exercises.  There are
a number of topics in common with the proposed book, although it includes much
less on message passing.  It does not include the following topics in the
proposed book: systematic testing; tuning; logging; encapsulation and
concurrent objects; linearization; fine-grained locking; sharding; conditions
(in the sense of the proposed book).

\emph{Principles of Concurrent and Distributed Programming}, M.~Ben-Ari,
Addison Wesley.\footnote{Note: this discussion is based on looking through
  the 1992 edition, and the table of contents of the 2005 edition.}  This is
another classic book, but is again rather dated.  I consider it broader but
shallower than the proposed book.  It covers semaphores,
monitors and message passing (each fairly briefly, with few examples).  It
does not seem to cover the following topics from the proposed book: systematic
testing; workers and controller pattern, including the bag-of-tasks pattern;
tuning; encapsulation; logging for debugging and testing; synchronisation
objects (beyond channels and barriers); different topologies for
communications; concurrent datatypes; linearization; synchronous parallel
programming; fine-grained locking; sharding.
%
Instead, it covers in more detail topics such as the critical section problem
(effectively building a lock from scratch), verification, the Linda
tuple-spaces paradigm, distributed algorithms, consensus, real-time systems.

\emph{Structured Parallel Programming}, Michael McCool, Arch D. Robison and
James Reinders, Morgan Kaufmann.
%
This book concentrates on ``high-level patterns'', by which they mean patterns
in specifications, rather than implementations.  They present ways of tackling
these patterns using various parallel libraries and compiler technologies.  In
particular, the book targets technologies produced by Intel (who employ all
three authors). My objection to this approach is that many problems fall
outside these patterns: programmers should know how to tackle problems
directly.  Further, the patterns considered are not difficult to tackle
directly: the proposed book includes examples of nearly all of them, and I
would expect readers to be able to implement the others for themselves.
%
The book seems to have negligible overlap with the proposed book: locks are
mentioned very briefly (less than three pages), as are barriers (less than one
page, with no examples); message-passing and semaphores seem not to be
covered.  None of the following topics is covered: systematic testing;
workers and controller pattern, including the bag-of-tasks pattern; tuning;
encapsulation; logging for debugging and testing; synchronisation objects;
different topologies for communications; concurrent datatypes; linearization;
sharding.
%
The writing style is very uneven: in some places it is over-dense, and seems
to assume a lot more programming experience than a typical student has; in
other places, it is rather plodding and repetitive.


%%%%%%%%%%%%%%%%%%%%%%%%%%%%%%%%%%%%%%%%%%%%%%%%%%%%%%%

\subsection*{Manuscript information}

The current draft is included with this proposal.  It is nearly complete, and
currently about 300 pages long (excluding Appendices~C and~D).  I would expect
the exercise solutions for students (included as Appendix C) to be available
online, and the solutions for instructors (included as Appendix D) to be
available from me.  The code would also be available online.

Chapters 1--9 are basically complete, with the exception of a few small items,
and maybe some extra exercises.  Chapters 10--12, the preface,  appendices,
and index need a more thorough review; in particular, a few more exercises
need adding.  The final chapter is currently little more than a first draft.
I think completion should take about three months.  In addition, the code for
online publication needs to be tidied up; this could be done in parallel with
editing.

%%%%%%%%%%%%%%%%%%%%%%%%%%%%%%%%%%%%%%%%%%%%%%%%%%%%%%%


\subsection*{Brief biography}

I am an Emeritus Professor of Computer Science at Oxford University, and an
Emeritus Fellow of St Catherine's College, Oxford.  My research has been in
the areas of Computer Security, Concurrency, Verification, and Concurrent
Programming.  I have over 30 years' teaching experience.  In particular, I
have lectured and tutored the subjects of this book, which has given me a very
good understanding of what students find difficult, and has helped me develop
ways of explaining the material.  I recently published the book \emph{Thinking
  About Programs} with ACM Press. 

\end{document}
